% --- IMPORTS ---
\documentclass[leqno]{article}
\usepackage{graphicx}
\usepackage{dsfont}
\usepackage{mathtools}
\usepackage{amssymb}
\usepackage{amsmath}
\usepackage{amsthm}
\usepackage{float}
\usepackage{ragged2e}
\usepackage{tensor}
\usepackage{fancyhdr}
\usepackage{empheq}
\usepackage[most]{tcolorbox}
\usepackage{enumitem}
\usepackage{dirtytalk}
\usepackage{marginnote}
\usepackage{microtype}
\usepackage[
  a4paper,
  left=0.8in,
  right=1in,
  top=1in,
  bottom=0.5in,
  headheight=14pt,
  headsep=0.4in
]{geometry}
\setlength{\parindent}{0cm}
% --- TikZ Packages ---
\usepackage{pgfplots}
\pgfplotsset{compat=1.18}
\usepgfplotslibrary{fillbetween, polar}
\usetikzlibrary{calc, positioning}
\usetikzlibrary{angles, quotes}
\usetikzlibrary{arrows.meta}
\usetikzlibrary{decorations.pathreplacing, decorations.markings}
\usetikzlibrary{patterns}
\usetikzlibrary{intersections}
% --- URL Packages ---
\usepackage[colorlinks=true,linkcolor=red,citecolor=magenta,urlcolor=black]{hyperref}%
\urlstyle{same}
% --- END OF IMPORTS ---

% --- CUSTOM COMMANDS ---
\RenewDocumentCommand{\marks}{o m}{%
  \IfValueTF{#1}%
    {\marginnote{\raggedright\textbf{[#2]}}[#1]}%
    {\marginnote{\raggedright\textbf{[#2]}}}%
}
\newcommand{\vspacerule}[2]{%
  \par\vspace*{#1}%
  \noindent\hrulefill\par
  \vspace*{#2}%
}
\definecolor{scarlet}{RGB}{205, 40, 75}
\newcommand{\cxl}[2]{%
  \tikz[baseline=(X.base)]{%
    \node[inner sep=0pt, outer sep=0pt, text=#1] (X) {$#2$};
    \draw[#1, line width=0.4pt] (X.south west)--(X.north east);
  }%
}
\newcommand{\questionitem}{%
  \item\phantomsection
  \addcontentsline{toc}{section}{Question \number\value{enumi}}%
}
\renewcommand{\contentsname}{Questions}
% --- END OF CUSTOM COMMANDS ---

% --- HEADER CONFIG ---
\pagestyle{fancy}
\fancyhead{}
\fancyfoot{}
\renewcommand{\headrulewidth}{0.9pt}
\lhead{\textit{Solutions} - Method of Differences}
\chead{}
\rhead{\href{https://danieljsmith.org}{danieljsmith.org}}
% --- END OF HEADER CONFIG ---

\begin{document}
\vspace*{20pt}
\tableofcontents
\newpage
\begin{enumerate}[
  label=\textbf{\arabic*.},
  leftmargin=*,
  topsep=0pt,
  partopsep=0pt,
  parsep=0pt
]

% ---- Question 1 ----
\questionitem
% [Source: _QBT___Solns__Series___Method_of_Differences, Q1]

\begin{enumerate}[label=\textbf{(\alph*)}]
  \item Show that, for $r>0$,
  \[
  \ln\left(1+\frac{2}{r(r+3)}\right)=\ln\left(\frac{r+1}{r}\right)-\ln\left(\frac{r+3}{r+2}\right)
  \]
  \marks[-30pt]{1}
  \item Hence, using the method of differences, show that for integer $n\geq 1$,
  \[
  \sum_{r=1}^{n}\ln\left(1+\frac{2}{r(r+3)}\right)=\ln\left(\frac{3(n+1)}{n+3}\right)
  \]
  \marks[-30pt]{4}
\end{enumerate}

\vspace*{10pt}

\begin{tcolorbox}[colback=white, colframe=scarlet, title={\textbf{Solution}}, breakable, grow to left by=6mm, grow to right by=10mm]
\begin{enumerate}[label=\textbf{(\alph*)}]
  \item Since \(r>0\), all logarithms are defined. Also,
\[
1+\frac{2}{r(r+3)}
=\frac{r(r+3)+2}{r(r+3)}
=\frac{r^2+3r+2}{r(r+3)}
=\frac{(r+1)(r+2)}{r(r+3)}
=\frac{r+1}{r}\cdot\frac{r+2}{r+3}
\]
So
\begin{align*}
\ln\left(1+\frac{2}{r(r+3)}\right)
&=\ln\left(\frac{r+1}{r}\right)+\ln\left(\frac{r+2}{r+3}\right) \\
&=\ln\left(\frac{r+1}{r}\right)-\ln\left(\frac{r+3}{r+2}\right)
\end{align*}
which is the required result.

  \item Let
\[
S_n=\sum_{r=1}^{n}\ln\left(1+\frac{2}{r(r+3)}\right)
\]
Using part (a),
\begin{align*}
S_n
&=\sum_{r=1}^{n}\left[\ln\left(\frac{r+1}{r}\right)-\ln\left(\frac{r+3}{r+2}\right)\right] \\
&=\left(\ln2-\ln\frac43\right)+\left(\ln\frac32-\ln\frac54\right)+\left(\ln\frac43-\ln\frac65\right)+\cdots+\left(\ln\frac{n+1}{n}-\ln\frac{n+3}{n+2}\right)
\end{align*}

Now write the positive and negative terms together:
\begin{align*}
S_n
&=\ln2+\ln\frac32+\cxl{blue}{\ln\frac43}+\cxl{red}{\ln\frac54}+\cxl{teal}{\ln\frac65}+\cdots+\cxl{violet}{\ln\frac{n+1}{n}} \\
&\quad-\cxl{blue}{\ln\frac43}-\cxl{red}{\ln\frac54}-\cxl{teal}{\ln\frac65}-\cdots-\cxl{violet}{\ln\frac{n+1}{n}}-\ln\frac{n+2}{n+1}-\ln\frac{n+3}{n+2}
\end{align*}

All the coloured middle terms cancel, so
\begin{align*}
S_n
&=\ln2+\ln\frac32-\ln\frac{n+2}{n+1}-\ln\frac{n+3}{n+2} \\
&=\ln\left(2\cdot\frac32\cdot\frac{n+1}{n+2}\cdot\frac{n+2}{n+3}\right) \\
&=\ln\left(\frac{3(n+1)}{n+3}\right)
\end{align*}

Therefore,
\[
\sum_{r=1}^{n}\ln\left(1+\frac{2}{r(r+3)}\right)=\ln\left(\frac{3(n+1)}{n+3}\right)
\]
\end{enumerate}
\end{tcolorbox}


\newpage

% ---- Question 2 ----
\questionitem
% [Source: _QBT___Solns__Series___Method_of_Differences, Q2]

Let $a$ be a positive constant. \\
\begin{enumerate}[label=\textbf{(\alph*)}]
  \item Use the method of differences to find
  \[
  \sum_{r=1}^{n} \frac{1}{(r+a)(r+a+1)(r+a+2)}
  \]
  in terms of $n$ and $a$. \marks{4} \\
  \item Hence find the value of $a$ for which
  \[
  \sum_{r=1}^{\infty} \frac{1}{(r+a)(r+a+1)(r+a+2)} = \frac{1}{24}
  \]
  \marks[-30pt]{3}
\end{enumerate}

\vspace*{10pt}

\begin{tcolorbox}[colback=white, colframe=scarlet, title={\textbf{Solution}}, breakable, grow to left by=6mm, grow to right by=10mm]
\begin{enumerate}[label=\textbf{(\alph*)}]
  \item Let
\[
S_n=\sum_{r=1}^{n}\frac{1}{(r+a)(r+a+1)(r+a+2)}
\]

To use the method of differences, first note that for any \(x\),
\begin{align*}
\frac{1}{x(x+1)}-\frac{1}{(x+1)(x+2)}
&=\frac{x+2-x}{x(x+1)(x+2)}\\
&=\frac{2}{x(x+1)(x+2)}
\end{align*}

So
\[
\frac{1}{x(x+1)(x+2)}
=\frac12\left(\frac{1}{x(x+1)}-\frac{1}{(x+1)(x+2)}\right)
\]

Now put \(x=r+a\). Then
\begin{align*}
S_n
&=\frac12\sum_{r=1}^{n}\left(\frac{1}{(r+a)(r+a+1)}-\frac{1}{(r+a+1)(r+a+2)}\right)
\end{align*}

Expanding the sum,
\begin{align*}
S_n
&=\frac12\Bigg[
\frac{1}{(a+1)(a+2)}
-\cxl{blue}{\frac{1}{(a+2)(a+3)}}\\
&\qquad
+\cxl{blue}{\frac{1}{(a+2)(a+3)}}
-\cxl{red}{\frac{1}{(a+3)(a+4)}}\\
&\qquad
+\cxl{red}{\frac{1}{(a+3)(a+4)}}
-\cdots\\
&\qquad
+\cxl{teal}{\frac{1}{(n+a)(n+a+1)}}
-\cxl{teal}{\frac{1}{(n+a)(n+a+1)}}
-\frac{1}{(n+a+1)(n+a+2)}
\Bigg]
\end{align*}

All the middle terms cancel, leaving
\[
S_n=\frac12\left(\frac{1}{(a+1)(a+2)}-\frac{1}{(n+a+1)(n+a+2)}\right)
\]

Hence
\[
\sum_{r=1}^{n}\frac{1}{(r+a)(r+a+1)(r+a+2)}
=\frac12\left(\frac{1}{(a+1)(a+2)}-\frac{1}{(n+a+1)(n+a+2)}\right)
\]

  \item From part (a),
\[
\sum_{r=1}^{n}\frac{1}{(r+a)(r+a+1)(r+a+2)}
=\frac12\left(\frac{1}{(a+1)(a+2)}-\frac{1}{(n+a+1)(n+a+2)}\right)
\]

As \(n\to\infty\),
\[
\frac{1}{(n+a+1)(n+a+2)}\to 0
\]

So
\[
\sum_{r=1}^{\infty}\frac{1}{(r+a)(r+a+1)(r+a+2)}
=\frac{1}{2(a+1)(a+2)}
\]

Given that this equals \(\frac{1}{24}\),
\begin{align*}
\frac{1}{2(a+1)(a+2)}&=\frac{1}{24}\\
2(a+1)(a+2)&=24\\
(a+1)(a+2)&=12\\
a^2+3a+2&=12\\
a^2+3a-10&=0\\
(a-2)(a+5)&=0
\end{align*}

So \(a=2\) or \(a=-5\). Since \(a\) is positive,
\[
a=2
\]

\end{enumerate}
\end{tcolorbox}


\newpage

% ---- Question 3 ----
\questionitem
% [Source: _QBT___Solns__Series___Method_of_Differences, Q3]

\begin{enumerate}[label=\textbf{(\alph*)}]
  \item Show that, for $r \ge 2$,
  \[
  \frac{4r}{\sqrt{r(r+2)}+\sqrt{r(r-2)}}=\sqrt{r(r+2)}-\sqrt{r(r-2)}
  \]
  \marks[-30pt]{2}
  \item Hence use the method of differences to determine
  \[
  \sum_{r=2}^{n} \frac{4r}{\sqrt{r(r+2)}+\sqrt{r(r-2)}}
  \]
  giving your answer in simplest form. \marks{3} \\
  \item Hence show that
  \[
  \sum_{r=3}^{7} \frac{4r}{\sqrt{r(r+2)}+\sqrt{r(r-2)}} = A\sqrt{2}+B\sqrt{3}+C\sqrt{7}
  \]
  where $A$, $B$ and $C$ are integers to be determined. \marks{2} \\
\end{enumerate}

\vspace*{10pt}

\begin{tcolorbox}[colback=white, colframe=scarlet, title={\textbf{Solution}}, breakable, grow to left by=6mm, grow to right by=10mm]
\begin{enumerate}[label=\textbf{(\alph*)}]
  \item Rationalise the denominator:
  \begin{align*}
  \frac{4r}{\sqrt{r(r+2)}+\sqrt{r(r-2)}}
  &=\frac{4r}{\sqrt{r(r+2)}+\sqrt{r(r-2)}}\cdot
  \frac{\sqrt{r(r+2)}-\sqrt{r(r-2)}}{\sqrt{r(r+2)}-\sqrt{r(r-2)}}\\
  &=\frac{4r\left(\sqrt{r(r+2)}-\sqrt{r(r-2)}\right)}{r(r+2)-r(r-2)}\\
  &=\frac{4r\left(\sqrt{r(r+2)}-\sqrt{r(r-2)}\right)}{4r}\\
  &=\sqrt{r(r+2)}-\sqrt{r(r-2)}
  \end{align*}
  Hence
  \[
  \frac{4r}{\sqrt{r(r+2)}+\sqrt{r(r-2)}}=\sqrt{r(r+2)}-\sqrt{r(r-2)}
  \]
  as required.

  \item Let
  \[
  u_r=\sqrt{r(r+2)}
  \]
  Then from part (a),
  \[
  \frac{4r}{\sqrt{r(r+2)}+\sqrt{r(r-2)}}=u_r-u_{r-2}
  \]
  So, if
  \[
  S_n=\sum_{r=2}^{n}\frac{4r}{\sqrt{r(r+2)}+\sqrt{r(r-2)}}
  \]
  then
  \begin{align*}
  S_n&=\sum_{r=2}^{n}(u_r-u_{r-2})\\
  &=\sum_{r=2}^{n}u_r-\sum_{r=2}^{n}u_{r-2}\\
  &=\sum_{r=2}^{n}u_r-\sum_{r=0}^{n-2}u_r
  \end{align*}
  using a change of index in the second sum.

  Writing out the terms shows the method of differences clearly:

  {
  \begin{align*}
  S_n&=(\cxl{blue}{u_2}+\cxl{red}{u_3}+\cxl{green!50!black}{u_4}+\cxl{orange}{u_5}+\cdots+\cxl{purple}{u_{n-2}}+u_{n-1}+u_n)\\
  &\qquad -(u_0+u_1+\cxl{blue}{u_2}+\cxl{red}{u_3}+\cxl{green!50!black}{u_4}+\cxl{orange}{u_5}+\cdots+\cxl{purple}{u_{n-2}})\\
  &=u_{n-1}+u_n-u_0-u_1
  \end{align*}
  }

  Now
  \[
  u_{n-1}=\sqrt{(n-1)(n+1)},\quad u_n=\sqrt{n(n+2)},\quad u_0=\sqrt{0\cdot 2}=0,\quad u_1=\sqrt{1\cdot 3}=\sqrt3
  \]
  Hence
  \[
  \sum_{r=2}^{n}\frac{4r}{\sqrt{r(r+2)}+\sqrt{r(r-2)}}
  =\sqrt{(n-1)(n+1)}+\sqrt{n(n+2)}-\sqrt3
  \]
  This is the required sum.

  \item Using part (b) with $n=7$,
  \begin{align*}
  \sum_{r=2}^{7}\frac{4r}{\sqrt{r(r+2)}+\sqrt{r(r-2)}}
  &=\sqrt{6\cdot 8}+\sqrt{7\cdot 9}-\sqrt3\\
  &=\sqrt{48}+\sqrt{63}-\sqrt3\\
  &=4\sqrt3+3\sqrt7-\sqrt3\\
  &=3\sqrt3+3\sqrt7
  \end{align*}

  The $r=2$ term is
  \[
  \frac{4(2)}{\sqrt{2(2+2)}+\sqrt{2(2-2)}}=\sqrt{2\cdot 4}-\sqrt{2\cdot 0}=\sqrt8=2\sqrt2
  \]
  from part (a).

  Therefore
  \begin{align*}
  \sum_{r=3}^{7}\frac{4r}{\sqrt{r(r+2)}+\sqrt{r(r-2)}}
  &=\sum_{r=2}^{7}\frac{4r}{\sqrt{r(r+2)}+\sqrt{r(r-2)}}-2\sqrt2\\
  &=3\sqrt3+3\sqrt7-2\sqrt2
  \end{align*}
  So
  \[
  \sum_{r=3}^{7}\frac{4r}{\sqrt{r(r+2)}+\sqrt{r(r-2)}}=-2\sqrt2+3\sqrt3+3\sqrt7
  \]
  Hence
  \[
  A=-2,\qquad B=3,\qquad C=3
  \]
\end{enumerate}
\end{tcolorbox}


\newpage

% ---- Question 4 ----
\questionitem
% [Source: _QBT___Solns__Series___Method_of_Differences, Q4]

Use the method of differences to show that
\[
\sum_{r=1}^{n} \frac{12}{(2r-1)(2r+1)(2r+3)} = \frac{n(an+b)}{(2n+1)(2n+3)}
\]
where $a$ and $b$ are integers to be determined. \marks{6} \\

\vspace*{10pt}

\begin{tcolorbox}[colback=white, colframe=scarlet, title={\textbf{Solution}}, breakable, grow to left by=6mm, grow to right by=10mm]
Let
\[
S_n=\sum_{r=1}^{n}\frac{12}{(2r-1)(2r+1)(2r+3)}
\]
Write the general term in partial fractions:
\[
\frac{12}{(2r-1)(2r+1)(2r+3)}
=\frac{A}{2r-1}+\frac{B}{2r+1}+\frac{C}{2r+3}
\]

Using the cover-up rule:
\[
A=\left.\frac{12}{(2r+1)(2r+3)}\right|_{\,2r-1=0}
=\frac{12}{2\cdot4}
=\frac32
\]
\[
B=\left.\frac{12}{(2r-1)(2r+3)}\right|_{\,2r+1=0}
=\frac{12}{(-2)(2)}
=-3
\]
\[
C=\left.\frac{12}{(2r-1)(2r+1)}\right|_{\,2r+3=0}
=\frac{12}{(-4)(-2)}
=\frac32
\]

So
\[
\frac{12}{(2r-1)(2r+1)(2r+3)}
=\frac{3}{2(2r-1)}-\frac{3}{2r+1}+\frac{3}{2(2r+3)}
\]

It is convenient to rewrite this as
\[
\frac{12}{(2r-1)(2r+1)(2r+3)}
=\frac32\left(\frac{1}{2r-1}-\frac{1}{2r+1}\right)
-\frac32\left(\frac{1}{2r+1}-\frac{1}{2r+3}\right)
\]

Hence
\begin{align*}
S_n
&=\sum_{r=1}^{n}\frac{12}{(2r-1)(2r+1)(2r+3)}\\
&=\frac32\sum_{r=1}^{n}\left(\frac{1}{2r-1}-\frac{1}{2r+1}\right)
-\frac32\sum_{r=1}^{n}\left(\frac{1}{2r+1}-\frac{1}{2r+3}\right)
\end{align*}

Expanding both telescoping sums:
\begin{align*}
S_n
&=\frac32\Bigg[
\left(1-\frac13\right)
+\left(\frac13-\frac15\right)
+\left(\frac15-\frac17\right)
+\cdots
+\left(\frac{1}{2n-1}-\frac{1}{2n+1}\right)
\Bigg]\\
&\qquad
-\frac32\Bigg[
\left(\frac13-\frac15\right)
+\left(\frac15-\frac17\right)
+\left(\frac17-\frac19\right)
+\cdots
+\left(\frac{1}{2n+1}-\frac{1}{2n+3}\right)
\Bigg]
\end{align*}

So
\begin{align*}
S_n
&=\frac32\Bigg[
1
-\cxl{blue}{\frac13}
+\cxl{blue}{\frac13}
-\cxl{red}{\frac15}
+\cxl{red}{\frac15}
-\cxl{green!60!black}{\frac17}
+\cxl{green!60!black}{\frac17}
-\cdots
-\frac{1}{2n+1}
\Bigg]\\
&\qquad
-\frac32\Bigg[
\frac13
-\cxl{red}{\frac15}
+\cxl{red}{\frac15}
-\cxl{green!60!black}{\frac17}
+\cxl{green!60!black}{\frac17}
-\cdots
-\frac{1}{2n+3}
\Bigg]
\end{align*}

Hence
\begin{align*}
S_n
&=\frac32\left(1-\frac{1}{2n+1}\right)
-\frac32\left(\frac13-\frac{1}{2n+3}\right)\\
&=\frac32-\frac{3}{2(2n+1)}-\frac12+\frac{3}{2(2n+3)}\\
&=1+\frac32\left(\frac{1}{2n+3}-\frac{1}{2n+1}\right)\\
&=1-\frac{3}{(2n+1)(2n+3)}\\
&=\frac{(2n+1)(2n+3)-3}{(2n+1)(2n+3)}\\
&=\frac{4n^2+8n}{(2n+1)(2n+3)}\\
&=\frac{n(4n+8)}{(2n+1)(2n+3)}
\end{align*}

Comparing with
\[
\frac{n(an+b)}{(2n+1)(2n+3)}
\]
gives
\[
a=4,\qquad b=8
\]

So
\[
\sum_{r=1}^{n}\frac{12}{(2r-1)(2r+1)(2r+3)}=\frac{n(4n+8)}{(2n+1)(2n+3)}
\]
and the integers are \(a=4\) and \(b=8\).
\end{tcolorbox}


\newpage

% ---- Question 5 ----
\questionitem
% [Source: _QBT___Solns__Series___Method_of_Differences, Q5]

\begin{enumerate}[label=\textbf{(\alph*)}]
  \item Use the method of differences to prove that for $n \ge 1$
  \[
  \sum_{r=1}^{n} \ln\left(1+\frac{2}{r(r+3)}\right)=\ln\left(\frac{3(n+1)}{n+3}\right)
  \]
  \marks[-30pt]{4}
  \item Hence find the exact value of
  \[
  \sum_{r=25}^{74} \ln\left(\left(1+\frac{2}{r(r+3)}\right)^{28}\right)
  \] \\
  Give your answer in the form $a\ln\left(\dfrac{b}{c}\right)$ where $a$, $b$ and $c$ are integers to be determined. \marks{3} \\
\end{enumerate}

\vspace*{10pt}

\begin{tcolorbox}[colback=white, colframe=scarlet, title={\textbf{Solution}}, breakable, grow to left by=6mm, grow to right by=10mm]
\begin{enumerate}[label=\textbf{(\alph*)}]
  \item Let
\[
S_n=\sum_{r=1}^{n}\ln\left(1+\frac{2}{r(r+3)}\right)
\]

First simplify the term inside the logarithm:
\begin{align*}
1+\frac{2}{r(r+3)}
&=\frac{r(r+3)+2}{r(r+3)}\\
&=\frac{r^2+3r+2}{r(r+3)}\\
&=\frac{(r+1)(r+2)}{r(r+3)}\\
&=\frac{r+1}{r}\cdot\frac{r+2}{r+3}
\end{align*}

So
\begin{align*}
\ln\left(1+\frac{2}{r(r+3)}\right)
&=\ln\left(\frac{r+1}{r}\cdot\frac{r+2}{r+3}\right)\\
&=\ln(r+1)-\ln r+\ln(r+2)-\ln(r+3)
\end{align*}

Hence
\begin{align*}
S_n
&=\sum_{r=1}^{n}\bigl(\ln(r+1)-\ln r\bigr)+\sum_{r=1}^{n}\bigl(\ln(r+2)-\ln(r+3)\bigr)
\end{align*}

Now expand each sum using the method of differences.

For the first sum,
\begin{align*}
\sum_{r=1}^{n}\bigl(\ln(r+1)-\ln r\bigr)
&=\bigl(\cxl{blue}{\ln2}+\cxl{red}{\ln3}+\cxl{teal}{\ln4}+\cdots+\cxl{violet}{\ln n}+\ln(n+1)\bigr)\\
&\qquad-\bigl(\ln1+\cxl{blue}{\ln2}+\cxl{red}{\ln3}+\cxl{teal}{\ln4}+\cdots+\cxl{violet}{\ln n}\bigr)\\
&=\ln(n+1)-\ln1
\end{align*}

For the second sum,
\begin{align*}
\sum_{r=1}^{n}\bigl(\ln(r+2)-\ln(r+3)\bigr)
&=\bigl(\ln3+\cxl{blue}{\ln4}+\cxl{red}{\ln5}+\cdots+\cxl{violet}{\ln(n+2)}\bigr)\\
&\qquad-\bigl(\cxl{blue}{\ln4}+\cxl{red}{\ln5}+\cdots+\cxl{violet}{\ln(n+2)}+\ln(n+3)\bigr)\\
&=\ln3-\ln(n+3)
\end{align*}

Therefore
\begin{align*}
S_n
&=\bigl(\ln(n+1)-\ln1\bigr)+\bigl(\ln3-\ln(n+3)\bigr)\\
&=\ln(n+1)+\ln3-\ln(n+3)
\end{align*}

Since $\ln1=0$,
\begin{align*}
S_n
&=\ln\left(\frac{3(n+1)}{n+3}\right)
\end{align*}

So
\[
\sum_{r=1}^{n}\ln\left(1+\frac{2}{r(r+3)}\right)=\ln\left(\frac{3(n+1)}{n+3}\right)
\]

  \item Using $\ln(a^{28})=28\ln a$,
\begin{align*}
\sum_{r=25}^{74}\ln\left(\left(1+\frac{2}{r(r+3)}\right)^{28}\right)
&=28\sum_{r=25}^{74}\ln\left(1+\frac{2}{r(r+3)}\right)
\end{align*}

From part (a), if
\[
S_n=\sum_{r=1}^{n}\ln\left(1+\frac{2}{r(r+3)}\right)
\]
then
\[
S_n=\ln\left(\frac{3(n+1)}{n+3}\right)
\]

So
\begin{align*}
\sum_{r=25}^{74}\ln\left(1+\frac{2}{r(r+3)}\right)
&=S_{74}-S_{24}\\
&=\ln\left(\frac{3(75)}{77}\right)-\ln\left(\frac{3(25)}{27}\right)\\
&=\ln\left(\frac{225}{77}\right)-\ln\left(\frac{75}{27}\right)\\
&=\ln\left(\frac{225}{77}\cdot\frac{27}{75}\right)\\
&=\ln\left(\frac{81}{77}\right)
\end{align*}

Hence
\begin{align*}
\sum_{r=25}^{74}\ln\left(\left(1+\frac{2}{r(r+3)}\right)^{28}\right)
&=28\ln\left(\frac{81}{77}\right)
\end{align*}

So the exact value is
\[
28\ln\left(\frac{81}{77}\right)
\]
\end{enumerate}
\end{tcolorbox}


\newpage

% ---- Question 6 ----
\questionitem
% [Source: _QBT___Solns__Series___Method_of_Differences, Q6]

Using partial fractions and a method of differences, prove that for positive integer $n$,
\[
\sum_{r=1}^{n} \frac{6}{r(r+1)(r+2)} = \frac{n(an+b)}{2(n+1)(n+2)}
\]
where $a$ and $b$ are constants to be found. \marks{6} \\

\vspace*{10pt}

\begin{tcolorbox}[colback=white, colframe=scarlet, title={\textbf{Solution}}, breakable, grow to left by=6mm, grow to right by=10mm]
Let
\[
S_n=\sum_{r=1}^{n}\frac{6}{r(r+1)(r+2)}
\]

Using partial fractions, write
\[
\frac{6}{r(r+1)(r+2)}=\frac{A}{r(r+1)}+\frac{B}{(r+1)(r+2)}
\]

Multiplying through by \(r(r+1)(r+2)\),
\[
6=A(r+2)+Br
\]

So
\[
6=(A+B)r+2A
\]

Comparing coefficients gives
\[
A+B=0,\qquad 2A=6
\]

Hence
\[
A=3,\qquad B=-3
\]

Therefore
\[
\frac{6}{r(r+1)(r+2)}=\frac{3}{r(r+1)}-\frac{3}{(r+1)(r+2)}
\]

So
\[
S_n=\sum_{r=1}^{n}\left(\frac{3}{r(r+1)}-\frac{3}{(r+1)(r+2)}\right)
\]

Now use the method of differences by writing out the first few and last few terms:
\begin{align*}
S_n
&=\left(\frac{3}{1\cdot2}-\cxl{blue}{\frac{3}{2\cdot3}}\right)
+\left(\cxl{blue}{\frac{3}{2\cdot3}}-\cxl{red}{\frac{3}{3\cdot4}}\right)
+\left(\cxl{red}{\frac{3}{3\cdot4}}-\cxl{teal}{\frac{3}{4\cdot5}}\right)
+\cdots \\
&\qquad
+\left(\cxl{violet}{\frac{3}{(n-1)n}}-\cxl{orange}{\frac{3}{n(n+1)}}\right)
+\left(\cxl{orange}{\frac{3}{n(n+1)}}-\frac{3}{(n+1)(n+2)}\right)
\end{align*}

All the coloured pairs cancel, leaving
\[
S_n=\frac{3}{1\cdot2}-\frac{3}{(n+1)(n+2)}
\]

So
\[
S_n=\frac{3}{2}-\frac{3}{(n+1)(n+2)}
\]

Put this over the common denominator \(2(n+1)(n+2)\):
\begin{align*}
S_n
&=\frac{3(n+1)(n+2)-6}{2(n+1)(n+2)} \\
&=\frac{3(n^2+3n+2)-6}{2(n+1)(n+2)} \\
&=\frac{3n^2+9n}{2(n+1)(n+2)} \\
&=\frac{3n(n+3)}{2(n+1)(n+2)}
\end{align*}

We are told this has the form
\[
\frac{n(an+b)}{2(n+1)(n+2)}
\]

Hence
\[
n(an+b)=3n(n+3)=n(3n+9)
\]

Since \(n\) is a positive integer, \(n\neq 0\), so
\[
an+b=3n+9
\]

Therefore
\[
a=3,\qquad b=9
\]

Hence
\[
\sum_{r=1}^{n}\frac{6}{r(r+1)(r+2)}=\frac{3n(n+3)}{2(n+1)(n+2)}
=\frac{n(3n+9)}{2(n+1)(n+2)}
\]
so the required constants are \(a=3\) and \(b=9\).
\end{tcolorbox}


\newpage

% ---- Question 7 ----
\questionitem
% [Source: _QBT___Solns__Series___Method_of_Differences, Q7]

Let $T_n = (n+1)^3 + (n+2)^3 + \cdots + (2n)^3$. \\
\begin{enumerate}[label=\textbf{(\alph*)}]
  \item By considering $(r+1)^4-(r-1)^4$, use the method of differences to prove that
  \[
  \sum_{r=1}^{n} r^3 = \frac{1}{4}n^2(n+1)^2
  \]
  \marks[-30pt]{5}
  \item Show that \[T_n = \frac{1}{4}n^2(an^2+bn+c)\] where $a$, $b$ and $c$ are integers to be determined. \marks{3} \\
  \item State the value of \[\lim_{n \to \infty} \frac{T_n}{n^4}\] \marks{1} \\
\end{enumerate}

\vspace*{10pt}

\begin{tcolorbox}[colback=white, colframe=scarlet, title={\textbf{Solution}}, breakable, grow to left by=6mm, grow to right by=10mm]
\begin{enumerate}[label=\textbf{(\alph*)}]
  \item First expand
\[
(r+1)^4-(r-1)^4
= \left(r^4+4r^3+6r^2+4r+1\right)-\left(r^4-4r^3+6r^2-4r+1\right)
=8r^3+8r
\]

Now sum from $r=1$ to $n$:
\[
\sum_{r=1}^n \left((r+1)^4-(r-1)^4\right)=8\sum_{r=1}^n r^3+8\sum_{r=1}^n r
\]

Using the method of differences on the left-hand side,
\begin{align*}
\sum_{r=1}^n \left((r+1)^4-(r-1)^4\right)
&=\left(\cxl{blue}{2^4}+\cxl{red}{3^4}+\cxl{teal}{4^4}+\cdots+\cxl{violet}{(n-1)^4}+n^4+(n+1)^4\right)\\
&\qquad-\left(0^4+1^4+\cxl{blue}{2^4}+\cxl{red}{3^4}+\cxl{teal}{4^4}+\cdots+\cxl{violet}{(n-1)^4}\right)\\
&=(n+1)^4+n^4-1
\end{align*}

Therefore
\[
(n+1)^4+n^4-1=8\sum_{r=1}^n r^3+8\sum_{r=1}^n r
\]

Using
\[
\sum_{r=1}^n r=\frac{n(n+1)}{2}
\]

we get
\begin{align*}
8\sum_{r=1}^n r^3
&=(n+1)^4+n^4-1-8\left(\frac{n(n+1)}{2}\right)\\
&=(n+1)^4+n^4-1-4n(n+1)\\
&=\left(n^4+4n^3+6n^2+4n+1\right)+n^4-1-4n^2-4n\\
&=2n^4+4n^3+2n^2\\
&=2n^2(n+1)^2
\end{align*}

So
\[
\sum_{r=1}^n r^3=\frac{2n^2(n+1)^2}{8}=\frac14 n^2(n+1)^2
\]

Hence
\[
\sum_{r=1}^{n} r^3=\frac14 n^2(n+1)^2
\]

  \item
\[
T_n=(n+1)^3+(n+2)^3+\cdots+(2n)^3=\sum_{r=n+1}^{2n} r^3
\]

So
\begin{align*}
T_n
&=\sum_{r=1}^{2n} r^3-\sum_{r=1}^{n} r^3\\
&=\frac14(2n)^2(2n+1)^2-\frac14 n^2(n+1)^2
\end{align*}

Now simplify:
\begin{align*}
T_n
&=\frac14\left(4n^2(2n+1)^2-n^2(n+1)^2\right)\\
&=\frac14\left(4n^2(4n^2+4n+1)-n^2(n^2+2n+1)\right)\\
&=\frac14\left(16n^4+16n^3+4n^2-n^4-2n^3-n^2\right)\\
&=\frac14\left(15n^4+14n^3+3n^2\right)\\
&=\frac14 n^2(15n^2+14n+3)
\end{align*}

Therefore
\[
T_n=\frac14 n^2(an^2+bn+c)
\quad\text{with}\quad
a=15,\ b=14,\ c=3
\]

  \item From part (b),
\[
\frac{T_n}{n^4}
=\frac14\left(15+\frac{14}{n}+\frac{3}{n^2}\right)
\]

As $n\to\infty$, $\dfrac{14}{n}\to 0$ and $\dfrac{3}{n^2}\to 0$, so
\[
\lim_{n\to\infty}\frac{T_n}{n^4}=\frac{15}{4}
\]
\end{enumerate}
\end{tcolorbox}


\newpage

% ---- Question 8 ----
\questionitem
% [Source: _QBT___Solns__Series___Method_of_Differences, Q8]

\begin{enumerate}[label=\textbf{(\alph*)}]
  \item By expressing $\dfrac{1}{r(r+2)}$ in partial fractions, show that
  \[
  \sum_{r=1}^{n} \frac{1}{r(r+2)} = \frac{3}{4} - \frac{1}{2(n+1)} - \frac{1}{2(n+2)}
  \]
  \marks[-30pt]{5}
  \item Hence determine the value of
  \[
  \sum_{r=1}^{\infty} \frac{1}{r(r+2)}
  \]
  \marks[-20pt]{2}
\end{enumerate}

\vspace*{10pt}

\begin{tcolorbox}[colback=white, colframe=scarlet, title={\textbf{Solution}}, breakable, grow to left by=6mm, grow to right by=10mm]
\begin{enumerate}[label=\textbf{(\alph*)}]
  \item Let
  \[
  S_n=\sum_{r=1}^{n}\frac{1}{r(r+2)}
  \]
  First express the general term in partial fractions:
  \[
  \frac{1}{r(r+2)}=\frac{A}{r}+\frac{B}{r+2}
  \]
  So
  \[
  1=A(r+2)+Br=(A+B)r+2A
  \]
  Comparing coefficients gives
  \[
  A+B=0,\qquad 2A=1
  \]
  Hence
  \[
  A=\frac12,\qquad B=-\frac12
  \]
  Therefore
  \[
  \frac{1}{r(r+2)}=\frac12\left(\frac1r-\frac1{r+2}\right)
  \]

  Substituting into the sum,
  \begin{align*}
  S_n
  &=\frac12\sum_{r=1}^{n}\left(\frac1r-\frac1{r+2}\right)\\
  &=\frac12\left[\left(1-\frac13\right)+\left(\frac12-\frac14\right)+\left(\frac13-\frac15\right)+\cdots+\left(\frac1n-\frac1{n+2}\right)\right]
  \end{align*}

  Now use the method of differences:
  \begin{align*}
  S_n
  &=\frac12\Bigl(1+\frac12+\cxl{blue}{\frac13}+\cxl{red}{\frac14}+\cdots+\cxl{teal}{\frac1n}
  -\cxl{blue}{\frac13}-\cxl{red}{\frac14}-\cdots-\cxl{teal}{\frac1n}
  -\frac1{n+1}-\frac1{n+2}\Bigr)\\
  &=\frac12\left(1+\frac12-\frac1{n+1}-\frac1{n+2}\right)\\
  &=\frac12\left(\frac32-\frac1{n+1}-\frac1{n+2}\right)\\
  &=\frac34-\frac1{2(n+1)}-\frac1{2(n+2)}
  \end{align*}

  Hence
  \[
  \sum_{r=1}^{n}\frac{1}{r(r+2)}=\frac34-\frac1{2(n+1)}-\frac1{2(n+2)}
  \]

  \item From part (a),
  \[
  \sum_{r=1}^{n}\frac{1}{r(r+2)}=\frac34-\frac1{2(n+1)}-\frac1{2(n+2)}
  \]
  Letting $n\to\infty$,
  \[
  \frac1{2(n+1)}\to 0,\qquad \frac1{2(n+2)}\to 0
  \]
  so
  \[
  \sum_{r=1}^{\infty}\frac{1}{r(r+2)}=\frac34
  \]

  The value of the infinite series is $\frac34$.
\end{enumerate}
\end{tcolorbox}


\newpage

% ---- Question 9 ----
\questionitem
% [Source: _QBT___Solns__Series___Method_of_Differences, Q9]

Show that
\[
\frac{1}{1 \times 4} + \frac{1}{2 \times 5} + \frac{1}{3 \times 6} + \cdots + \frac{1}{n(n+3)} = \frac{11}{18} - \frac{an^2+bn+c}{3(n+1)(n+2)(n+3)}
\]
where $a$, $b$ and $c$ are constants to be found. \marks{6} \\

\vspace*{10pt}

\begin{tcolorbox}[colback=white, colframe=scarlet, title={\textbf{Solution}}, breakable, grow to left by=6mm, grow to right by=10mm]
Let
\[
S=\frac{1}{1\times 4}+\frac{1}{2\times 5}+\frac{1}{3\times 6}+\cdots+\frac{1}{n(n+3)}
=\sum_{r=1}^n \frac{1}{r(r+3)}
\]

Use partial fractions on the general term:
\[
\frac{1}{r(r+3)}=\frac{A}{r}+\frac{B}{r+3}
\]

So
\[
1=A(r+3)+Br=(A+B)r+3A
\]

Comparing coefficients:
\[
A+B=0,\qquad 3A=1
\]

Hence
\[
A=\frac13,\qquad B=-\frac13
\]

Therefore
\[
\frac{1}{r(r+3)}=\frac13\left(\frac1r-\frac1{r+3}\right)
\]

So
\[
S=\frac13\sum_{r=1}^n\left(\frac1r-\frac1{r+3}\right)
\]

Now write out the terms to use the method of differences:
\begin{align*}
S
&=\frac13\Biggl[
1+\frac12+\frac13+\cxl{blue}{\frac14}+\cxl{red}{\frac15}+\cdots+\cxl{teal}{\frac1n}\\
&\qquad\qquad\qquad
-\cxl{blue}{\frac14}-\cxl{red}{\frac15}-\cdots-\cxl{teal}{\frac1n}-\frac1{n+1}-\frac1{n+2}-\frac1{n+3}
\Biggr]
\end{align*}

Everything from $\frac14$ to $\frac1n$ cancels, leaving
\[
S=\frac13\left(1+\frac12+\frac13-\frac1{n+1}-\frac1{n+2}-\frac1{n+3}\right)
\]

Now
\[
1+\frac12+\frac13=\frac{6+3+2}{6}=\frac{11}{6}
\]

Hence
\[
S=\frac{11}{18}-\frac13\left(\frac1{n+1}+\frac1{n+2}+\frac1{n+3}\right)
\]

Combine the three fractions:
\[
\frac1{n+1}+\frac1{n+2}+\frac1{n+3}
=\frac{(n+2)(n+3)+(n+1)(n+3)+(n+1)(n+2)}{(n+1)(n+2)(n+3)}
\]

Expand the numerator:
\begin{align*}
(n+2)(n+3)&=n^2+5n+6 \\
(n+1)(n+3)&=n^2+4n+3 \\
(n+1)(n+2)&=n^2+3n+2
\end{align*}

So
\[
(n+2)(n+3)+(n+1)(n+3)+(n+1)(n+2)=3n^2+12n+11
\]

Therefore
\[
S=\frac{11}{18}-\frac{3n^2+12n+11}{3(n+1)(n+2)(n+3)}
\]

Comparing with
\[
\frac{11}{18}-\frac{an^2+bn+c}{3(n+1)(n+2)(n+3)}
\]
we get
\[
a=3,\qquad b=12,\qquad c=11
\]

So the required result is
\[
\frac{1}{1 \times 4} + \frac{1}{2 \times 5} + \frac{1}{3 \times 6} + \cdots + \frac{1}{n(n+3)}
= \frac{11}{18} - \frac{3n^2+12n+11}{3(n+1)(n+2)(n+3)}
\]
\end{tcolorbox}


\newpage

% ---- Question 10 ----
\questionitem
% [Source: _QBT___Solns__Series___Method_of_Differences, Q10]

\begin{enumerate}[label=\textbf{(\alph*)}]
  \item Use standard results from the formula booklet to show that
  \[
  \sum_{r=1}^{N} r(2r-1)(2r+1)=\frac{1}{2}N(N+1)(2N^2+2N-1)
  \]
  \marks[-30pt]{3}
  \item Express $\dfrac{6r+13}{(2r-1)(2r+3)}$ in partial fractions and hence use the method of differences to find
  \[
  \sum_{r=1}^{N} \frac{6r+13}{(2r-1)(2r+3)}\left(\frac{1}{2}\right)^{r+1}
  \]
  in terms of $N$. \marks{4} \\
  \item Deduce the value of
  \[
  \sum_{r=1}^{\infty} \frac{6r+13}{(2r-1)(2r+3)}\left(\frac{1}{2}\right)^{r+1}
  \]
  \marks[-20pt]{1}
\end{enumerate}

\vspace*{10pt}

\begin{tcolorbox}[colback=white, colframe=scarlet, title={\textbf{Solution}}, breakable, grow to left by=6mm, grow to right by=10mm]
\begin{enumerate}[label=\textbf{(\alph*)}]
  \item Using
\[
(2r-1)(2r+1)=4r^2-1
\]
we get
\[
r(2r-1)(2r+1)=r(4r^2-1)=4r^3-r
\]

So
\[
\sum_{r=1}^{N} r(2r-1)(2r+1)=4\sum_{r=1}^{N} r^3-\sum_{r=1}^{N} r
\]

From the formula book,
\[
\sum_{r=1}^{N} r=\frac{N(N+1)}{2},
\qquad
\sum_{r=1}^{N} r^3=\left(\frac{N(N+1)}{2}\right)^2
\]

Hence
\begin{align*}
\sum_{r=1}^{N} r(2r-1)(2r+1)
&=4\left(\frac{N(N+1)}{2}\right)^2-\frac{N(N+1)}{2} \\
&=N^2(N+1)^2-\frac12 N(N+1) \\
&=\frac12 N(N+1)\bigl(2N(N+1)-1\bigr) \\
&=\frac12 N(N+1)(2N^2+2N-1)
\end{align*}

So
\[
\sum_{r=1}^{N} r(2r-1)(2r+1)=\frac12 N(N+1)(2N^2+2N-1)
\]

  \item Let
\[
\frac{6r+13}{(2r-1)(2r+3)}=\frac{A}{2r-1}+\frac{B}{2r+3}
\]

Then
\[
6r+13=A(2r+3)+B(2r-1)
\]
so
\[
6r+13=(2A+2B)r+(3A-B)
\]

Comparing coefficients,
\[
2A+2B=6,\qquad 3A-B=13
\]
so
\[
A+B=3,\qquad 3A-B=13
\]

Adding gives \(4A=16\), so \(A=4\), and then \(B=-1\).

Therefore
\[
\frac{6r+13}{(2r-1)(2r+3)}=\frac{4}{2r-1}-\frac{1}{2r+3}
\]

Now let
\[
S_N=\sum_{r=1}^{N} \frac{6r+13}{(2r-1)(2r+3)}\left(\frac12\right)^{r+1}
\]

Then
\begin{align*}
S_N
&=\sum_{r=1}^{N}\left(\frac{4}{2r-1}-\frac{1}{2r+3}\right)\frac{1}{2^{r+1}} \\
&=\sum_{r=1}^{N}\left(\frac{4}{(2r-1)2^{r+1}}-\frac{1}{(2r+3)2^{r+1}}\right) \\
&=\sum_{r=1}^{N}\left(\frac{4}{(2r-1)2^{r+1}}-\frac{4}{(2r+3)2^{r+3}}\right)
\end{align*}

Define
\[
u_r=\frac{1}{(2r-1)2^{r+1}}
\]
Then
\[
u_{r+2}=\frac{1}{(2r+3)2^{r+3}}
\]
so
\[
S_N=4\sum_{r=1}^{N}(u_r-u_{r+2})
\]

Using the method of differences,
\begin{align*}
S_N
&=4\Bigl[(u_1+u_2+\cxl{blue}{u_3}+\cxl{red}{u_4}+\cdots+\cxl{teal}{u_N})\\
&\qquad\qquad-(\cxl{blue}{u_3}+\cxl{red}{u_4}+\cdots+\cxl{teal}{u_N}+u_{N+1}+u_{N+2})\Bigr] \\
&=4(u_1+u_2-u_{N+1}-u_{N+2})
\end{align*}

Now
\[
u_1=\frac{1}{1\cdot 2^2}=\frac14,
\qquad
u_2=\frac{1}{3\cdot 2^3}=\frac{1}{24}
\]
and
\[
u_{N+1}=\frac{1}{(2N+1)2^{N+2}},
\qquad
u_{N+2}=\frac{1}{(2N+3)2^{N+3}}
\]

So
\begin{align*}
S_N
&=4\left(\frac14+\frac{1}{24}-\frac{1}{(2N+1)2^{N+2}}-\frac{1}{(2N+3)2^{N+3}}\right) \\
&=\frac76-\frac{1}{(2N+1)2^N}-\frac{1}{(2N+3)2^{N+1}}
\end{align*}

Combining the two final fractions,
\begin{align*}
\frac{1}{(2N+1)2^N}+\frac{1}{(2N+3)2^{N+1}}
&=\frac{2}{(2N+1)2^{N+1}}+\frac{1}{(2N+3)2^{N+1}} \\
&=\frac{2(2N+3)+(2N+1)}{(2N+1)(2N+3)2^{N+1}} \\
&=\frac{6N+7}{(2N+1)(2N+3)2^{N+1}}
\end{align*}

Hence
\[
\sum_{r=1}^{N} \frac{6r+13}{(2r-1)(2r+3)}\left(\frac12\right)^{r+1}
=\frac76-\frac{1}{(2N+1)2^N}-\frac{1}{(2N+3)2^{N+1}}
\]
or equivalently
\[
\sum_{r=1}^{N} \frac{6r+13}{(2r-1)(2r+3)}\left(\frac12\right)^{r+1}
=\frac76-\frac{6N+7}{(2N+1)(2N+3)2^{N+1}}
\]

  \item From part (b),
\[
S_N=\frac76-\frac{6N+7}{(2N+1)(2N+3)2^{N+1}}
\]

As \(N\to\infty\), the second term tends to \(0\) because the denominator contains the factor \(2^{N+1}\).

Therefore
\[
\sum_{r=1}^{\infty} \frac{6r+13}{(2r-1)(2r+3)}\left(\frac12\right)^{r+1}
=\frac76
\]
\end{enumerate}
\end{tcolorbox}


\newpage

% ---- Question 11 ----
\questionitem
% [Source: _QBT___Solns__Series___Method_of_Differences, Q11]

\begin{enumerate}[label=\textbf{(\alph*)}]
  \item Show that
  \[
  (2n)^5-(2n-2)^5 \equiv 10(2n-1)^4+20(2n-1)^2+2
  \]
  \marks[-20pt]{2}
  \item Hence, using the method of differences, show that for all positive integer values of $n$,
  \[
  \sum_{r=1}^{n}(2r-1)^4=\frac{1}{15}n(2n-1)(2n+1)(an^2+bn+c)
  \]
  where $a$, $b$ and $c$ are integers to be determined. \marks{7} \\
\end{enumerate}

\vspace*{10pt}

\begin{tcolorbox}[colback=white, colframe=scarlet, title={\textbf{Solution}}, breakable, grow to left by=6mm, grow to right by=10mm]
\begin{enumerate}[label=\textbf{(\alph*)}]
  \item Let \(m=2n-1\). Then
\[
2n=m+1,\qquad 2n-2=m-1
\]
so
\begin{align*}
(2n)^5-(2n-2)^5
&=(m+1)^5-(m-1)^5\\
&=\left(m^5+5m^4+10m^3+10m^2+5m+1\right)\\
&\qquad-\left(m^5-5m^4+10m^3-10m^2+5m-1\right)\\
&=10m^4+20m^2+2
\end{align*}
Substituting back \(m=2n-1\),
\[
(2n)^5-(2n-2)^5=10(2n-1)^4+20(2n-1)^2+2
\]
So the identity is shown.

  \item Let
\[
S=\sum_{r=1}^{n}(2r-1)^4
\]
Using the result from part (a) with \(n\) replaced by \(r\),
\[
(2r)^5-(2r-2)^5=10(2r-1)^4+20(2r-1)^2+2
\]
Now sum from \(r=1\) to \(n\):
\[
\sum_{r=1}^{n}\left((2r)^5-(2r-2)^5\right)
=
10\sum_{r=1}^{n}(2r-1)^4+20\sum_{r=1}^{n}(2r-1)^2+2n
\]
That is,
\[
\sum_{r=1}^{n}\left((2r)^5-(2r-2)^5\right)=10S+20\sum_{r=1}^{n}(2r-1)^2+2n
\]

For the left-hand side, use the method of differences:
\begin{align*}
\sum_{r=1}^{n}\left((2r)^5-(2r-2)^5\right)
&=\Bigl(\cxl{blue}{2^5}+\cxl{red}{4^5}+\cxl{teal}{6^5}+\cdots+\cxl{violet}{(2n-2)^5}+(2n)^5\Bigr)\\
&\qquad-\Bigl(0^5+\cxl{blue}{2^5}+\cxl{red}{4^5}+\cxl{teal}{6^5}+\cdots+\cxl{violet}{(2n-2)^5}\Bigr)\\
&=(2n)^5\\
&=32n^5
\end{align*}
So
\[
32n^5=10S+20\sum_{r=1}^{n}(2r-1)^2+2n
\]

Now find \(\sum_{r=1}^{n}(2r-1)^2\):
\begin{align*}
\sum_{r=1}^{n}(2r-1)^2
&=\sum_{r=1}^{n}(4r^2-4r+1)\\
&=4\sum_{r=1}^{n}r^2-4\sum_{r=1}^{n}r+\sum_{r=1}^{n}1\\
&=4\cdot \frac{n(n+1)(2n+1)}{6}-4\cdot \frac{n(n+1)}{2}+n\\
&=\frac{2n(n+1)(2n+1)}{3}-2n(n+1)+n\\
&=\frac{2n(n+1)(2n+1)-6n(n+1)+3n}{3}\\
&=\frac{n\left(2(n+1)(2n+1)-6(n+1)+3\right)}{3}\\
&=\frac{n(4n^2-1)}{3}\\
&=\frac{n(2n-1)(2n+1)}{3}
\end{align*}

Substitute this into the equation for \(S\):
\begin{align*}
32n^5
&=10S+20\left(\frac{n(2n-1)(2n+1)}{3}\right)+2n
\end{align*}
Hence
\begin{align*}
10S
&=32n^5-\frac{20n(2n-1)(2n+1)}{3}-2n\\
&=32n^5-\frac{20n(4n^2-1)}{3}-2n\\
&=32n^5-\frac{80}{3}n^3+\frac{20}{3}n-2n\\
&=32n^5-\frac{80}{3}n^3+\frac{14}{3}n
\end{align*}
So
\begin{align*}
S
&=\frac{1}{10}\left(32n^5-\frac{80}{3}n^3+\frac{14}{3}n\right)\\
&=\frac{48n^5-40n^3+7n}{15}
\end{align*}

Now factorise:
\begin{align*}
48n^5-40n^3+7n
&=n(48n^4-40n^2+7)\\
&=n(4n^2-1)(12n^2-7)\\
&=n(2n-1)(2n+1)(12n^2-7)
\end{align*}
Therefore
\[
\sum_{r=1}^{n}(2r-1)^4=\frac{1}{15}n(2n-1)(2n+1)(12n^2-7)
\]

Comparing with
\[
\sum_{r=1}^{n}(2r-1)^4=\frac{1}{15}n(2n-1)(2n+1)(an^2+bn+c)
\]
gives
\[
a=12,\qquad b=0,\qquad c=-7
\]
\end{enumerate}
\end{tcolorbox}


\newpage

% ---- Question 12 ----
\questionitem
% [Source: _QBT___Solns__Series___Method_of_Differences, Q12]

Let
\[
S_n = \sum_{r=1}^{n} \frac{1}{(3r-2)(3r+4)}
\]
where $n \geq 1$. \\

\begin{enumerate}[label=\textbf{(\alph*)}]
  \item Use the method of differences to show that
  \[
  S_n = \frac{n(15n+17)}{8(3n+1)(3n+4)}
  \]
  \marks[-30pt]{6}
  \item Show that, for any number $k$ greater than $\dfrac{24}{5}$, if the difference between $\dfrac{5}{24}$ and $S_n$ is less than $\dfrac{1}{k}$, then \\[5pt]
  \[
  n > \frac{k - 15 + \sqrt{k^2 + 81}}{18}
  \]
  \marks[-30pt]{6}
\end{enumerate}

\vspace*{10pt}

\begin{tcolorbox}[colback=white, colframe=scarlet, title={\textbf{Solution}}, breakable, grow to left by=6mm, grow to right by=10mm]
\begin{enumerate}[label=\textbf{(\alph*)}]
  \item Write
\[
\frac{1}{(3r-2)(3r+4)}=\frac{A}{3r-2}+\frac{B}{3r+4}
\]

Then
\[
1=A(3r+4)+B(3r-2)
\]

Comparing coefficients of \(r\) and constants gives
\[
3A+3B=0,\qquad 4A-2B=1
\]

So
\[
A+B=0 \implies B=-A
\]
and hence
\[
4A-2(-A)=1 \implies 6A=1 \implies A=\frac16,\quad B=-\frac16
\]

Therefore
\[
\frac{1}{(3r-2)(3r+4)}=\frac16\left(\frac1{3r-2}-\frac1{3r+4}\right)
\]

So
\[
S_n=\frac16\sum_{r=1}^{n}\left(\frac1{3r-2}-\frac1{3r+4}\right)
\]

Now expand the sum:
\[
\begin{aligned}
S_n
&=\frac16\Bigg[\left(1-\frac17\right)+\left(\frac14-\frac1{10}\right)+\left(\frac17-\frac1{13}\right)+\left(\frac1{10}-\frac1{16}\right)+\cdots \\
&\hspace{3cm}+\left(\frac1{3n-5}-\frac1{3n+1}\right)+\left(\frac1{3n-2}-\frac1{3n+4}\right)\Bigg]
\end{aligned}
\]

Rearranging to show the cancellation pattern,
\[
\begin{aligned}
S_n
&=\frac16\Bigg(1+\frac14+\cxl{blue}{\frac17}+\cxl{red}{\frac1{10}}+\cdots+\cxl{teal}{\frac1{3n-2}} \\
&\hspace{2.3cm}-\cxl{blue}{\frac17}-\cxl{red}{\frac1{10}}-\cdots-\cxl{teal}{\frac1{3n-2}}-\frac1{3n+1}-\frac1{3n+4}\Bigg)
\end{aligned}
\]

The coloured middle terms cancel, leaving
\[
S_n=\frac16\left(1+\frac14-\frac1{3n+1}-\frac1{3n+4}\right)
\]

Now simplify:
\[
S_n=\frac16\left(\frac54-\frac{(3n+4)+(3n+1)}{(3n+1)(3n+4)}\right)
\]
\[
S_n=\frac16\left(\frac54-\frac{6n+5}{(3n+1)(3n+4)}\right)
\]

Putting the terms inside the bracket over a common denominator,
\[
\begin{aligned}
S_n
&=\frac16\left(\frac{5(3n+1)(3n+4)-4(6n+5)}{4(3n+1)(3n+4)}\right) \\
&=\frac16\left(\frac{5(9n^2+15n+4)-24n-20}{4(3n+1)(3n+4)}\right) \\
&=\frac16\left(\frac{45n^2+75n+20-24n-20}{4(3n+1)(3n+4)}\right) \\
&=\frac16\left(\frac{45n^2+51n}{4(3n+1)(3n+4)}\right) \\
&=\frac16\left(\frac{3n(15n+17)}{4(3n+1)(3n+4)}\right)
\end{aligned}
\]

Hence
\[
S_n=\frac{n(15n+17)}{8(3n+1)(3n+4)}
\]

  \item From part (a),
\[
S_n=\frac16\left(\frac54-\frac{6n+5}{(3n+1)(3n+4)}\right)
=\frac5{24}-\frac{6n+5}{6(3n+1)(3n+4)}
\]

So
\[
\frac5{24}-S_n=\frac{6n+5}{6(3n+1)(3n+4)}
\]

If this difference is less than \(\frac1k\), then
\[
\frac{6n+5}{6(3n+1)(3n+4)}<\frac1k
\]

Since \(k>\frac{24}{5}>0\), all quantities are positive, so we may cross-multiply:
\[
k(6n+5)<6(3n+1)(3n+4)
\]

Expanding:
\[
6kn+5k<6(9n^2+15n+4)=54n^2+90n+24
\]

Hence
\[
54n^2+(90-6k)n+(24-5k)>0
\]

Now solve the corresponding quadratic equation
\[
54n^2+(90-6k)n+(24-5k)=0
\]

Using the quadratic formula,
\[
n=\frac{-(90-6k)\pm\sqrt{(90-6k)^2-4(54)(24-5k)}}{108}
\]

Simplify the discriminant:
\[
\begin{aligned}
(90-6k)^2-4(54)(24-5k)
&=36(15-k)^2-216(24-5k) \\
&=36\left((15-k)^2-6(24-5k)\right) \\
&=36\left(k^2-30k+225-144+30k\right) \\
&=36(k^2+81)
\end{aligned}
\]

Therefore
\[
\begin{aligned}
n
&=\frac{6k-90\pm 6\sqrt{k^2+81}}{108} \\
&=\frac{k-15\pm\sqrt{k^2+81}}{18}
\end{aligned}
\]

Since the coefficient of \(n^2\) is positive, the inequality
\[
54n^2+(90-6k)n+(24-5k)>0
\]
holds for values of \(n\) outside the two roots.

Also,
\[
k-15-\sqrt{k^2+81}<k-15-k=-15<0
\]

So the smaller root is negative. As \(n\geq 1\), it follows that
\[
n>\frac{k-15+\sqrt{k^2+81}}{18}
\]

This is the required result.
\end{enumerate}
\end{tcolorbox}


\end{enumerate}
\end{document}